% =============================================================================
% МЕТОДИЧКА: ДИЗАЙН СХВАТОК И РАЗМЕЩЕНИЕ ПРОТИВНИКОВ
% Для Climax Arena с двухчастной структурой
% =============================================================================
\documentclass[a4paper,11pt]{report}

\usepackage[utf8]{inputenc}
\usepackage[T2A]{fontenc}
\usepackage[russian]{babel}

\usepackage[left=2.5cm, right=2.5cm, top=2.5cm, bottom=2.5cm, headheight=14pt]{geometry}

\usepackage{tikz}
\usetikzlibrary{
  shapes.geometric, arrows.meta, positioning, calc,
  patterns, decorations.pathreplacing, backgrounds,
  fit, shadows, matrix, angles, quotes
}
\usepackage{pgfplots}
\pgfplotsset{compat=1.18}
\usepackage{graphicx}
\graphicspath{{images/}}

\usepackage{microtype}
\usepackage{enumitem}
\usepackage{tabularx}
\usepackage{booktabs}
\usepackage{multirow}
\usepackage{float}
\usepackage{wrapfig}
\usepackage{amssymb}

% --- Colors ---
\usepackage{xcolor}
\definecolor{titanblue}{HTML}{1A5276}
\definecolor{titanred}{HTML}{C0392B}
\definecolor{titanorange}{HTML}{E67E22}
\definecolor{titangreen}{HTML}{27AE60}
\definecolor{titandark}{HTML}{2C3E50}
\definecolor{titanlight}{HTML}{ECF0F1}
\definecolor{titangold}{HTML}{F39C12}
\definecolor{titancyan}{HTML}{1ABC9C}
\definecolor{polaritypos}{HTML}{3498DB}
\definecolor{polarityneg}{HTML}{E74C3C}

% --- Headers ---
\usepackage{fancyhdr}
\pagestyle{fancy}
\fancyhf{}
\fancyhead[L]{\small\textcolor{titandark}{\leftmark}}
\fancyhead[R]{\small\textcolor{titandark}{Encounter Design Guide}}
\fancyfoot[C]{\thepage}
\renewcommand{\headrulewidth}{0.4pt}

% --- Boxes ---
\usepackage[most]{tcolorbox}

\newtcolorbox{principle}[1]{
  colback=titanblue!5, colframe=titanblue,
  fonttitle=\bfseries, title={#1},
  arc=2mm, boxrule=0.8pt,
  left=6pt, right=6pt, top=4pt, bottom=4pt
}
\newtcolorbox{warning}[1]{
  colback=titanred!5, colframe=titanred,
  fonttitle=\bfseries, title={#1},
  arc=2mm, boxrule=0.8pt,
  left=6pt, right=6pt, top=4pt, bottom=4pt
}
\newtcolorbox{example}[1]{
  colback=titangreen!5, colframe=titangreen,
  fonttitle=\bfseries, title={#1},
  arc=2mm, boxrule=0.8pt,
  left=6pt, right=6pt, top=4pt, bottom=4pt
}
\newtcolorbox{tip}[1]{
  colback=titanorange!5, colframe=titanorange,
  fonttitle=\bfseries, title={#1},
  arc=2mm, boxrule=0.8pt,
  left=6pt, right=6pt, top=4pt, bottom=4pt
}
\newtcolorbox{polarity}[1]{
  colback=titancyan!5, colframe=titancyan,
  fonttitle=\bfseries, title={#1},
  arc=2mm, boxrule=0.8pt,
  left=6pt, right=6pt, top=4pt, bottom=4pt
}
\newtcolorbox{rulebox}[1]{
  colback=titangold!5, colframe=titangold,
  fonttitle=\bfseries, title={#1},
  arc=2mm, boxrule=0.8pt,
  left=6pt, right=6pt, top=4pt, bottom=4pt
}

\usepackage{hyperref}
\hypersetup{
  colorlinks=true,
  linkcolor=titanblue,
  urlcolor=titanblue,
}

\usepackage{titlesec}
\titleformat{\chapter}[display]
  {\normalfont\huge\bfseries\color{titandark}}
  {\chaptertitlename\ \thechapter}{20pt}{\Huge}
\titleformat{\section}
  {\normalfont\Large\bfseries\color{titanblue}}
  {\thesection}{1em}{}
\titleformat{\subsection}
  {\normalfont\large\bfseries\color{titandark}}
  {\thesubsection}{1em}{}
\titleformat{\subsubsection}
  {\normalfont\normalsize\bfseries\color{titandark}}
  {\thesubsubsection}{1em}{}

\setcounter{tocdepth}{3}
\setcounter{secnumdepth}{3}

% =============================================================================
\begin{document}

% --- Title Page ---
\begin{titlepage}
\begin{tikzpicture}[remember picture, overlay]
  \fill[titandark] (current page.south west) rectangle (current page.north east);
  \foreach \i in {1,...,15} {
    \pgfmathsetmacro{\x}{rand*21}
    \pgfmathsetmacro{\y}{rand*29.7}
    \fill[titanred!20, opacity=0.12] (\x,\y) circle ({0.4+rand*0.6});
  }
  \fill[titanred] (0,16) rectangle (21,16.15);
  \fill[titanorange] (0,15.85) rectangle (21,16);
  \node[anchor=south, text=white, font=\fontsize{30}{36}\selectfont\bfseries]
    at (10.5,18) {ENCOUNTER DESIGN};
  \node[anchor=south, text=titanred, font=\fontsize{24}{30}\selectfont\bfseries]
    at (10.5,17) {GUIDE};
  \node[anchor=north, text=titanlight, font=\Large, text width=14cm, align=center]
    at (10.5,15.5) {%
      Размещение противников и дизайн схваток\\[4pt]
      для двухчастной Climax Arena\\[12pt]
      \textcolor{titanorange}{Часть I:} Коридорная подготовка (5--10 мин)\\
      \textcolor{titanred}{Часть II:} Центральная арена 100$\times$100$\times$50м (6--9 мин)
    };
  \node[anchor=south, text=titanlight!60, font=\small]
    at (10.5,2) {Polarity Project --- Internal Document --- 2026};
\end{tikzpicture}
\end{titlepage}

\tableofcontents
\newpage


% =============================================================================
\chapter{Архитектура двухчастной Climax Arena}
% =============================================================================

\section{Общая структура}

Climax Arena состоит из двух принципиально разных зон,
каждая из которых требует своего подхода к размещению врагов:

\begin{tikzpicture}[scale=0.85]
  % Part 1: Building
  \draw[very thick, titanblue, rounded corners=4pt, fill=titanblue!5]
    (0,0) rectangle (6,8);
  \node[font=\Large\bfseries, titanblue] at (3,7.3) {ЧАСТЬ I};
  \node[font=\small\bfseries, titandark] at (3,6.6) {Подготовка};
  \node[font=\small, titandark] at (3,6) {5--10 минут};

  % Building floors
  \draw[thick, titandark!30] (0.5,1) -- (5.5,1);
  \draw[thick, titandark!30] (0.5,2.5) -- (5.5,2.5);
  \draw[thick, titandark!30] (0.5,4) -- (5.5,4);
  \node[font=\tiny, titandark!50] at (3,0.5) {Этаж 1};
  \node[font=\tiny, titandark!50] at (3,1.7) {Этаж 2};
  \node[font=\tiny, titandark!50] at (3,3.2) {Этаж 3};
  \node[font=\tiny, titandark!50] at (3,4.7) {Этаж 4};

  % Circular path
  \draw[very thick, titanorange, -{Stealth}, dashed]
    (1,0.5) -- (5,0.5) -- (5,1.7) -- (1,1.7) -- (1,3.2) -- (5,3.2) -- (5,4.7) -- (3,5.5);

  % Enemy groups
  \fill[titanred] (4,0.5) circle (0.15);
  \fill[titanred] (4.3,0.5) circle (0.15);
  \fill[titanred] (2,1.7) circle (0.15);
  \fill[titanred] (2.3,1.7) circle (0.15);
  \fill[titanred] (2.6,1.7) circle (0.15);
  \fill[titanred] (4,3.2) circle (0.15);
  \fill[titanred] (4.3,3.2) circle (0.15);
  \fill[titanred] (4.6,3.2) circle (0.15);
  \fill[titanred] (4.9,3.2) circle (0.15);

  % Arrow to part 2
  \draw[ultra thick, titanred, -{Stealth[length=10pt]}] (6.5,4) -- (8,4);
  \node[font=\small\bfseries, titanred, above] at (7.25,4) {Переход};

  % Part 2: Arena
  \draw[very thick, titanred, rounded corners=4pt, fill=titanred!5]
    (8.5,0) rectangle (15.5,8);
  \node[font=\Large\bfseries, titanred] at (12,7.3) {ЧАСТЬ II};
  \node[font=\small\bfseries, titandark] at (12,6.6) {Центральная битва};
  \node[font=\small, titandark] at (12,6) {6--9 минут};

  % Arena interior
  \draw[thick, titandark!20, fill=titandark!3] (9.5,1) rectangle (14.5,5.5);
  \node[font=\scriptsize, titandark!40] at (12,5) {100m $\times$ 100m $\times$ 50m};

  % Enemy swarm
  \foreach \x/\y in {10.2/2, 10.8/3.5, 11.5/1.5, 12/4, 12.5/2.5,
    13/3.8, 13.5/1.8, 11/4.5, 13.8/4.2, 10.5/4.8} {
    \fill[titanred] (\x,\y) circle (0.12);
  }

  % Player path
  \draw[very thick, titangreen, dashed, -{Stealth}]
    (9.8,3) to[out=30,in=180] (11,4.5) to[out=0,in=150]
    (13,3) to[out=-30,in=180] (14,2);

  % Bottom labels
  \node[font=\small, text=titandark, text width=5cm, align=center] at (3,-1) {
    \textbf{Малые группы}\\
    2--4 противника\\
    Коридоры + комнаты\\
    Линейный прогресс
  };
  \node[font=\small, text=titandark, text width=5cm, align=center] at (12,-1) {
    \textbf{Большие группы}\\
    8--20 противников\\
    Открытое пространство\\
    3D-движение (полёт)
  };
\end{tikzpicture}

\section{Ключевые различия между частями}

\begin{table}[H]
\centering
\caption{Сравнение двух частей Climax Arena}
\begin{tabularx}{\textwidth}{@{}lXX@{}}
\toprule
\textbf{Параметр} & \textbf{Часть I: Подготовка} & \textbf{Часть II: Центральная битва} \\
\midrule
Пространство & Коридоры, комнаты, лестницы & Открытая арена 100$\times$100$\times$50м \\
Движение & Wallrun, slide, double jump & Полный 3D-полёт \\
Группы врагов & 2--4 противника & 8--20+ одновременно \\
Стычки & 4--6 отдельных энкаунтеров & 3--4 волны + босс \\
Темп & Staccato (бой-пауза-бой) & Crescendo (непрерывный рост) \\
Давление & Фронтальное, фланговое & Сферическое (360\textdegree{} + вертикаль) \\
Навык & Позиционирование, аккуратность & Полный skill expression \\
\bottomrule
\end{tabularx}
\end{table}


% =============================================================================
\chapter{Роли противников}
% =============================================================================

\section{Система ролей}

Каждый тип врага выполняет конкретную \textbf{функцию} в бою.
Функция определяет, \textbf{какое давление} враг создаёт на игрока
и \textbf{какую реакцию} вынуждает.

\begin{principle}{Принцип DOOM Eternal (Hugo Martin)}
Враги~--- это <<шахматные фигуры>>, которые двигают игрока по арене.
Каждый тип врага создаёт уникальное давление, заставляя игрока
менять позицию, оружие или тактику.
\end{principle}

\subsection{Пять основных ролей}

\begin{tikzpicture}[
  role/.style={
    rectangle, rounded corners=4pt,
    minimum width=3.8cm, minimum height=2.2cm,
    draw=#1, fill=#1!8, text=titandark,
    font=\small, align=center
  }
]
  % Roles
  \node[role=titanred] (rush) at (0,6) {
    \textbf{\color{titanred}RUSHER}\\[2pt]
    {\scriptsize Ближний бой}\\
    {\scriptsize Быстрый, слабый}\\
    {\scriptsize Давит на \textbf{позицию}}
  };
  \node[role=titanblue] (anchor) at (5,6) {
    \textbf{\color{titanblue}ANCHOR}\\[2pt]
    {\scriptsize Средняя/дальняя дист.}\\
    {\scriptsize Стреляет с позиции}\\
    {\scriptsize Давит на \textbf{укрытие}}
  };
  \node[role=titanorange] (disrupt) at (10,6) {
    \textbf{\color{titanorange}DISRUPTOR}\\[2pt]
    {\scriptsize Гранаты, AoE, щиты}\\
    {\scriptsize Выбивает из укрытия}\\
    {\scriptsize Давит на \textbf{движение}}
  };
  \node[role=titangold] (tank) at (2.5,2) {
    \textbf{\color{titangold}TANK}\\[2pt]
    {\scriptsize Медленный, живучий}\\
    {\scriptsize Огромный урон вблизи}\\
    {\scriptsize Давит на \textbf{дистанцию}}
  };
  \node[role=titancyan] (spec) at (7.5,2) {
    \textbf{\color{titancyan}SPECIALIST}\\[2pt]
    {\scriptsize Уникальная механика}\\
    {\scriptsize Снайпер, целитель, баффер}\\
    {\scriptsize Давит на \textbf{приоритеты}}
  };

  % Pressure arrows
  \draw[very thick, titanred!50, -{Stealth}] (rush.south) -- ++(0,-0.5)
    node[below, font=\tiny] {$\to$ Игрок убегает};
  \draw[very thick, titanblue!50, -{Stealth}] (anchor.south) -- ++(0,-0.5)
    node[below, font=\tiny] {$\to$ Игрок прячется};
  \draw[very thick, titanorange!50, -{Stealth}] (disrupt.south) -- ++(0,-0.5)
    node[below, font=\tiny] {$\to$ Игрок перемещается};
  \draw[very thick, titangold!50, -{Stealth}] (tank.south) -- ++(0,-0.5)
    node[below, font=\tiny] {$\to$ Игрок отступает};
  \draw[very thick, titancyan!50, -{Stealth}] (spec.south) -- ++(0,-0.5)
    node[below, font=\tiny] {$\to$ Игрок фокусирует};
\end{tikzpicture}

\subsection{Таблица ролей с параметрами}

\begin{table}[H]
\centering
\caption{Параметры ролей противников}
\begin{tabularx}{\textwidth}{@{}l c c c c X@{}}
\toprule
\textbf{Роль} & \textbf{HP} & \textbf{Скорость} & \textbf{Урон} & \textbf{Дист.} & \textbf{Реакция игрока} \\
\midrule
Rusher & Низкий & Высокая & Средний & Ближняя & Отстрел на подходе или уклонение \\
Anchor & Средний & Низкая & Средний & Средняя & Поиск укрытия, фланг \\
Disruptor & Средний & Средняя & Низкий & Любая & Смена позиции, приоритизация \\
Tank & Высокий & Низкая & Высокий & Ближняя & Кайтинг, дистанция \\
Specialist & Низкий & Разная & Разный & Дальняя & Немедленное уничтожение \\
\bottomrule
\end{tabularx}
\end{table}

\section{Комбинации ролей: давление как система}

\begin{principle}{Принцип эмерджентности (F.E.A.R.)}
Когда два типа врагов работают вместе, тактики, эффективные
против каждого по отдельности, \textbf{перестают работать}.
Это вынуждает игрока искать новые решения.
\end{principle}

\begin{table}[H]
\centering
\caption{Комбинации ролей и возникающее давление}
\begin{tabularx}{\textwidth}{@{}l l X@{}}
\toprule
\textbf{Комбинация} & \textbf{Конфликт} & \textbf{Что происходит} \\
\midrule
Rusher + Anchor & Бежать vs прятаться & Нельзя стоять (rusher) и нельзя бежать (anchor стреляет). Нужно \textbf{маневрировать} \\
Anchor + Disruptor & Прятаться vs двигаться & Нельзя сидеть в укрытии (disruptor выбивает), нельзя выходить (anchor стреляет). Нужна \textbf{агрессия} \\
Rusher + Tank & Убегать vs отступать & Rusher догоняет, Tank блокирует путь отступления. Нужно \textbf{приоритизировать} \\
Tank + Specialist & Дистанция vs фокус & Хочется держать дистанцию (tank), но specialist нужно убить первым. \textbf{Риск/награда} \\
Rusher + Specialist & Позиция vs приоритет & Rusher не даёт спокойно целиться в specialist-а. Нужны \textbf{быстрые решения} \\
\bottomrule
\end{tabularx}
\end{table}

\begin{warning}{Антипаттерн: Однородные группы}
Группа из 6 одинаковых Rusher-ов~--- скучно. Игрок применяет
одну тактику (стрелять в подбегающих) и не принимает решений.
Группа из 3 Rusher + 1 Anchor~--- уже интересно: убивать rusher-ов
или сначала нейтрализовать anchor?
\end{warning}

\begin{rulebox}{Правило <<Нет безопасной тактики>>}
Хорошая комбинация врагов не имеет \textbf{единственного правильного
ответа}. Если игрок может каждый раз делать одно и то же (например,
всегда прыгать на высоту и стрелять вниз)~--- комбинация неудачная.
Добавьте врага, который \textbf{наказывает} эту тактику.
\end{rulebox}


% =============================================================================
\chapter{Часть I: Коридорные стычки}
% =============================================================================

\section{Контекст}

Игрок бегает по большому зданию, перемещаясь между этажами по
кругу, продвигаясь к входу в центральную арену. Стычки с малыми
группами (2--4 врага) происходят в коридорах, на лестницах,
в комнатах и на открытых балконах.

\section{Правила размещения врагов в коридорах}

\subsection{Правило 1: Дистанция обнаружения}

\begin{tikzpicture}[scale=0.75]
  \node[font=\bfseries\color{titandark}] at (7,6.5) {Зоны дистанции в коридоре};

  % Corridor
  \fill[titanlight] (0,2) rectangle (14,4);
  \draw[thick] (0,2) -- (14,2);
  \draw[thick] (0,4) -- (14,4);

  % Player
  \fill[titangreen] (1,3) circle (0.3);
  \node[font=\scriptsize\bfseries, titangreen, below] at (1,1.7) {PLAYER};

  % Distance zones
  \fill[titangreen!15] (1,2) rectangle (4,4);
  \fill[titanorange!15] (4,2) rectangle (8,4);
  \fill[titanred!15] (8,2) rectangle (14,4);

  % Zone labels
  \node[font=\scriptsize\bfseries, titangreen] at (2.5,5) {Мёртвая зона};
  \node[font=\tiny, titangreen] at (2.5,4.5) {0--5м};
  \node[font=\tiny, titangreen] at (2.5,1.5) {Слишком близко};
  \node[font=\tiny, titangreen] at (2.5,1) {для спавна};

  \node[font=\scriptsize\bfseries, titanorange] at (6,5) {Зона реакции};
  \node[font=\tiny, titanorange] at (6,4.5) {5--15м};
  \node[font=\tiny, titanorange] at (6,1.5) {Игрок видит и};
  \node[font=\tiny, titanorange] at (6,1) {может среагировать};

  \node[font=\scriptsize\bfseries, titanred] at (11,5) {Зона инициации};
  \node[font=\tiny, titanred] at (11,4.5) {15--30м};
  \node[font=\tiny, titanred] at (11,1.5) {Идеально для};
  \node[font=\tiny, titanred] at (11,1) {размещения группы};

  % Enemy group in initiation zone
  \fill[titanred] (10,3) circle (0.2);
  \fill[titanred] (10.5,2.7) circle (0.2);
  \fill[titanred] (11,3.3) circle (0.2);

  % Dashed lines
  \draw[thick, dashed, titandark!40] (4,1.5) -- (4,5.5);
  \draw[thick, dashed, titandark!40] (8,1.5) -- (8,5.5);
\end{tikzpicture}

\begin{rulebox}{Дистанция обнаружения}
\begin{itemize}[leftmargin=*]
  \item \textbf{0--5м (мёртвая зона):} Никогда не спавните врагов прямо
    перед носом игрока. Исключение: засады, но только если враг
    стоит \textbf{спиной} к игроку
  \item \textbf{5--15м (зона реакции):} Игрок видит врага и может
    среагировать. Подходит для одиночных Rusher-ов, которые бегут
    навстречу~--- игрок успевает сориентироваться
  \item \textbf{15--30м (зона инициации):} Идеальная дистанция для
    размещения группы. Игрок видит состав, оценивает угрозу, выбирает
    тактику. Это та дистанция, на которой бой \textbf{начинается осмысленно}
  \item \textbf{30м+ (зона за поворотом):} Враги не видны~--- размещаются
    за углом, за дверью, на следующем этаже. Создают ощущение
    <<населённости>> здания
\end{itemize}
\end{rulebox}

\subsection{Правило 2: Привязка к геометрии}

\begin{principle}{Принцип F.E.A.R.}
Задача дизайнера~--- создать пространство с интересными элементами
(укрытия, окна, фланговые проходы, столы), которые AI \textbf{эксплуатирует}.
Враги не запрограммированы на фланг специально~--- они ищут укрытие
ближе к игроку, и фланговое поведение возникает \textbf{естественно}.
\end{principle}

Каждая группа врагов должна быть привязана к конкретному
\textbf{геометрическому элементу}, который делает бой интересным:

\begin{table}[H]
\centering
\caption{Привязка врагов к геометрии}
\begin{tabularx}{\textwidth}{@{}lXX@{}}
\toprule
\textbf{Геометрия} & \textbf{Какие враги} & \textbf{Почему интересно} \\
\midrule
Т-образный перекрёсток & 2 Anchor по сторонам & Игрок не может стрелять в обоих одновременно \\
Лестничный пролёт & 2 Rusher сверху & Используют высоту как преимущество \\
Балкон над коридором & 1 Anchor + 1 Rusher снизу & Вертикальное давление + горизонтальное \\
Комната с 2 входами & 3 Anchor внутри & Игрок выбирает, с какого входа атаковать \\
Длинный коридор & 1 Specialist в конце & Дальняя угроза, нужно закрыть дистанцию \\
Комната с колоннами & 2 Rusher + 1 Disruptor & Колонны дают укрытие, но disruptor выбивает \\
\bottomrule
\end{tabularx}
\end{table}

\subsection{Правило 3: Направление давления}

\begin{tikzpicture}[scale=0.75]
  \node[font=\bfseries\color{titandark}] at (7,8) {Три паттерна давления в коридорах};

  % Pattern 1: Frontal
  \fill[titanlight] (0,5) rectangle (4.3,7);
  \draw[thick] (0,5) -- (4.3,5) (0,7) -- (4.3,7);
  \fill[titangreen] (0.5,6) circle (0.25);
  \fill[titanred] (3.5,5.7) circle (0.15);
  \fill[titanred] (3.5,6.3) circle (0.15);
  \draw[thick, titanred, -{Stealth}] (3.3,5.7) -- (1.3,5.9);
  \draw[thick, titanred, -{Stealth}] (3.3,6.3) -- (1.3,6.1);
  \node[font=\scriptsize\bfseries, titandark] at (2.15,7.5) {A: Фронтальное};
  \node[font=\tiny, titandark] at (2.15,4.5) {Простое, предсказуемое};

  % Pattern 2: Crossfire
  \fill[titanlight] (5,5) rectangle (9.3,7);
  \draw[thick] (5,5) -- (9.3,5) (5,7) -- (9.3,7);
  \fill[titangreen] (5.5,6) circle (0.25);
  \fill[titanred] (8.5,6.8) circle (0.15);
  \fill[titanred] (8.5,5.2) circle (0.15);
  \draw[thick, titanred, -{Stealth}] (8.3,6.7) -- (6.3,6.2);
  \draw[thick, titanred, -{Stealth}] (8.3,5.3) -- (6.3,5.8);
  \node[font=\scriptsize\bfseries, titandark] at (7.15,7.5) {B: Перекрёстный};
  \node[font=\tiny, titandark] at (7.15,4.5) {Требует приоритизации};

  % Pattern 3: Sandwich
  \fill[titanlight] (10,5) rectangle (14.3,7);
  \draw[thick] (10,5) -- (14.3,5) (10,7) -- (14.3,7);
  \fill[titangreen] (12.15,6) circle (0.25);
  \fill[titanred] (10.5,6) circle (0.15);
  \fill[titanred] (13.8,6) circle (0.15);
  \fill[titanred] (13.5,5.5) circle (0.15);
  \draw[thick, titanred, -{Stealth}] (10.7,6) -- (11.7,6);
  \draw[thick, titanred, -{Stealth}] (13.6,6) -- (12.6,6);
  \node[font=\scriptsize\bfseries, titandark] at (12.15,7.5) {C: Тиски};
  \node[font=\tiny, titandark] at (12.15,4.5) {Максимальное давление};

  % Legend
  \node[font=\small, text=titandark, text width=14cm, align=left] at (7,3) {
    \textbf{A) Фронтальное:} все враги впереди. Используйте для первых
    стычек~--- знакомство.\\
    \textbf{B) Перекрёстное:} враги на разных высотах или за разными
    укрытиями. Игрок не может спрятаться от всех сразу.\\
    \textbf{C) Тиски (Sandwich):} враги с двух сторон. Используйте
    \textbf{максимум один раз} за подготовку~--- очень стрессовый паттерн.
  };
\end{tikzpicture}


\section{Эскалация стычек по мере продвижения}

По мере подъёма по зданию стычки должны \textbf{эскалировать}
по четырём осям:

\begin{tikzpicture}
  \begin{axis}[
    width=14cm, height=6cm,
    xlabel={Этаж / прогресс},
    ylabel={Значение},
    xmin=0, xmax=5,
    ymin=0, ymax=5,
    xtick={1,2,3,4},
    xticklabels={Этаж 1, Этаж 2, Этаж 3, Этаж 4},
    legend style={at={(0.02,0.98)}, anchor=north west, font=\scriptsize},
    grid=major,
    grid style={titanlight},
    axis lines=left,
    title={\textbf{Эскалация параметров по этажам}},
    title style={font=\bfseries\color{titandark}}
  ]
    \addplot[very thick, titanred, mark=*] coordinates {(1,1) (2,2) (3,3) (4,4)};
    \addlegendentry{Кол-во врагов}
    \addplot[very thick, titanblue, mark=square*] coordinates {(1,1) (2,1) (3,2) (4,3)};
    \addlegendentry{Типов ролей}
    \addplot[very thick, titanorange, mark=triangle*] coordinates {(1,1) (2,2) (3,2) (4,3)};
    \addlegendentry{Направлений давления}
    \addplot[very thick, titangreen, mark=diamond*] coordinates {(1,1) (2,1) (3,2) (4,2)};
    \addlegendentry{Вертикальность}
  \end{axis}
\end{tikzpicture}

\subsection{Детальный план стычек}

\begin{table}[H]
\centering
\caption{План стычек для четырёхэтажной подготовки}
\begin{tabularx}{\textwidth}{@{}c c l X c@{}}
\toprule
\textbf{Этаж} & \textbf{Стычка} & \textbf{Состав} & \textbf{Геометрия и тактика} & \textbf{Давление} \\
\midrule
\multirow{2}{*}{1} & 1a & 2 Rusher & Коридор, фронтальные, бегут навстречу & Фронт \\
 & 1b & 2 Anchor & Комната, стреляют из-за укрытий & Фронт \\
\midrule
\multirow{2}{*}{2} & 2a & 2R + 1A & Т-перекрёсток: anchor в конце, rusher-ы с боков & Перекрёст \\
 & 2b & 3 Anchor & Балкон: стреляют снизу, игрок сверху & Вертик. \\
\midrule
\multirow{2}{*}{3} & 3a & 2R + 1D & Комната с колоннами: disruptor выбивает из-за колонн & Перекрёст \\
 & 3b & 2A + 1Sp & Длинный зал: specialist в конце, anchor-ы в середине & Вертик. \\
\midrule
4 & 4 & 2R+1A+1D & Финальная комната перед ареной: все роли & Тиски \\
\bottomrule
\end{tabularx}
\end{table}

\begin{tip}{Между стычками}
Между каждой стычкой должен быть \textbf{коридор без врагов}
длиной 5--15 секунд бега. Это зона отдыха: подобрать ресурсы,
осмотреться, заметить окно с видом на арену (foreshadowing).
\end{tip}

\section{Spawn-логика в коридорах}

\begin{rulebox}{5 правил спавна в Части I}
\begin{enumerate}[leftmargin=*]
  \item \textbf{Не спавнить за спиной.} В линейном здании враги за
    спиной~--- нечестно. Исключение: scripted засада (тиски),
    максимум 1 раз
  \item \textbf{Спавнить до прибытия игрока.} Враги уже должны \textbf{быть}
    в комнате, когда игрок входит. Появление из воздуха~--- последнее средство
  \item \textbf{Использовать spawn closets.} Двери, лифты, вентиляция~---
    логичные точки, из которых враги <<приходят>>
  \item \textbf{Аудио-фидбек.} Игрок должен \textbf{слышать} врагов до того,
    как увидит (шаги, голоса, звуки оружия). Это даёт время подготовиться
  \item \textbf{Одна стычка = одно пространство.} Не смешивайте врагов
    из разных комнат. Каждая стычка~--- изолированный энкаунтер
    с чётким началом и концом
\end{enumerate}
\end{rulebox}


% =============================================================================
\chapter{Часть II: Центральная арена (100$\times$100$\times$50м)}
% =============================================================================

\section{Принципы боя в большом объёме}

Центральная арена~--- это не комната, а \textbf{трёхмерное боевое
пространство}. Игрок активно летает, используя все механики.
Это радикально меняет правила размещения врагов.

\begin{principle}{Принцип сферического давления}
В 3D-арене враг может быть \textbf{сверху, снизу, сбоку, сзади}.
Укрытия теряют ценность. Единственная защита~--- \textbf{движение}.
Враги размещаются так, чтобы не существовало <<безопасной высоты>>
или <<безопасного угла>>.
\end{principle}

\subsection{Зоны арены}

\begin{tikzpicture}[scale=0.65]
  \node[font=\bfseries\color{titandark}] at (7,12) {Вертикальный разрез арены (50м)};

  % Arena outline
  \draw[very thick, titandark] (0,0) rectangle (14,10);

  % Vertical zones
  \fill[titanblue!10] (0,0) rectangle (14,3);
  \fill[titanorange!8] (0,3) rectangle (14,7);
  \fill[titanred!8] (0,7) rectangle (14,10);

  \draw[thick, dashed, titandark!30] (0,3) -- (14,3);
  \draw[thick, dashed, titandark!30] (0,7) -- (14,7);

  % Zone labels
  \node[font=\bfseries, titanblue] at (7,1.5) {GROUND ZONE (0--15м)};
  \node[font=\scriptsize, titandark] at (7,0.7) {Tank-и, тяжёлые враги, ресурсы};

  \node[font=\bfseries, titanorange] at (7,5) {MID ZONE (15--35м)};
  \node[font=\scriptsize, titandark] at (7,4.3) {Основной бой, anchor-ы, платформы};

  \node[font=\bfseries, titanred] at (7,8.5) {HIGH ZONE (35--50м)};
  \node[font=\scriptsize, titandark] at (7,7.7) {Specialist-ы, снайперы, ключевые цели};

  % Platforms
  \fill[titandark!20] (1,3) rectangle (3.5,3.3);
  \fill[titandark!20] (5,5) rectangle (7.5,5.3);
  \fill[titandark!20] (9,4) rectangle (11,4.3);
  \fill[titandark!20] (2,7) rectangle (4,7.3);
  \fill[titandark!20] (10,7.5) rectangle (12.5,7.8);

  % Enemies in zones
  \fill[titangold] (2,0.8) circle (0.3);
  \node[font=\tiny, white] at (2,0.8) {T};
  \fill[titanred] (6,5.7) circle (0.2);
  \fill[titanred] (10,4.8) circle (0.2);
  \fill[titancyan] (3,7.8) circle (0.2);
  \fill[titancyan] (11,8.5) circle (0.2);

  % Player
  \fill[titangreen] (7,6.5) circle (0.3);
  \node[font=\tiny, white] at (7,6.5) {P};
  \draw[thick, titangreen, dashed, -{Stealth}] (7,6.5) to[out=30,in=180] (9.5,7);
  \draw[thick, titangreen, dashed, -{Stealth}] (7,6.5) to[out=-120,in=90] (5,4.5);

  % Side annotation
  \node[font=\small, text=titandark, text width=4cm, align=left, right] at (14.5,5) {
    Игрок свободно\\
    перемещается\\
    между зонами.\\[4pt]
    Враги привязаны\\
    к своим зонам\\
    и создают\\
    давление оттуда.
  };
\end{tikzpicture}

\subsection{Горизонтальные зоны}

\begin{tikzpicture}[scale=0.65]
  \node[font=\bfseries\color{titandark}] at (7,11.5) {Вид сверху: горизонтальные зоны (100$\times$100м)};

  % Arena outline
  \draw[very thick, titandark] (0,0) rectangle (14,10);

  % Concentric zones
  \draw[thick, dashed, titanred!50] (3,2.5) rectangle (11,7.5);
  \draw[thick, dashed, titanorange!50] (1.5,1.2) rectangle (12.5,8.8);

  % Zone labels
  \fill[titanred!8] (3,2.5) rectangle (11,7.5);
  \node[font=\bfseries, titanred] at (7,5) {KILL ZONE};
  \node[font=\tiny, titandark] at (7,4.3) {Максимальная концентрация врагов};
  \node[font=\tiny, titandark] at (7,3.8) {и ресурсов. Центр арены};

  \node[font=\scriptsize\bfseries, titanorange, rotate=90] at (1,5) {FLANK ZONE};
  \node[font=\scriptsize\bfseries, titanorange, rotate=-90] at (13,5) {FLANK ZONE};

  \node[font=\scriptsize\bfseries, titanblue] at (7,0.5) {EDGE ZONE};
  \node[font=\scriptsize\bfseries, titanblue] at (7,9.5) {EDGE ZONE};
  \node[font=\tiny, titanblue] at (7,10.3) {Перезарядка, ресурсы, передышка};

  % Resource pickups at edges
  \fill[titangreen] (0.5,5) circle (0.2);
  \fill[titangreen] (13.5,5) circle (0.2);
  \fill[titangreen] (7,0.3) circle (0.2);
  \fill[titangreen] (7,9.7) circle (0.2);
  \node[font=\tiny, titangreen] at (0.5,5.5) {HP};
  \node[font=\tiny, titangreen] at (13.5,5.5) {HP};

  % Enemy positions (example wave)
  \foreach \x/\y in {5/4, 6/6, 8/5, 9/3.5, 7/6.5, 5.5/5, 8.5/6} {
    \fill[titanred] (\x,\y) circle (0.15);
  }
\end{tikzpicture}

\begin{rulebox}{Зоны арены}
\begin{itemize}[leftmargin=*]
  \item \textbf{Kill Zone} (центр, ~60$\times$60м): Основная зона боя.
    Здесь размещается большинство врагов. Вход сюда = бой
  \item \textbf{Flank Zone} (периферия, ~20м шириной): Менее
    плотная зона. Отдельные враги на платформах. Путь для
    перегруппировки
  \item \textbf{Edge Zone} (стены, ~10м): Относительная безопасность.
    Ресурсы (HP, боеприпасы) размещены здесь, чтобы \textbf{вознаградить}
    игрока за рискованный рывок из центра
\end{itemize}
\end{rulebox}

\section{Дизайн волн}

\subsection{Паттерны волновой эскалации}

\begin{tikzpicture}[scale=0.85]
  \begin{axis}[
    width=14cm, height=6.5cm,
    xlabel={Время (минуты)},
    ylabel={Кол-во врагов на арене},
    xmin=0, xmax=9,
    ymin=0, ymax=22,
    xtick={0,1,...,9},
    grid=major,
    grid style={titanlight},
    axis lines=left,
    title={\textbf{Волновая эскалация центральной арены}},
    title style={font=\bfseries\color{titandark}},
    legend style={at={(0.02,0.98)}, anchor=north west, font=\scriptsize}
  ]
    % Wave boundaries
    \draw[thick, dashed, titandark!30] (axis cs:2.5,0) -- (axis cs:2.5,22);
    \draw[thick, dashed, titandark!30] (axis cs:5,0) -- (axis cs:5,22);
    \draw[thick, dashed, titandark!30] (axis cs:7,0) -- (axis cs:7,22);

    \node[font=\scriptsize\bfseries, titanblue] at (axis cs:1.25,21) {Волна 1};
    \node[font=\scriptsize\bfseries, titanorange] at (axis cs:3.75,21) {Волна 2};
    \node[font=\scriptsize\bfseries, titanred] at (axis cs:6,21) {Волна 3};
    \node[font=\scriptsize\bfseries, titangold] at (axis cs:8,21) {BOSS};

    % Enemy count curve
    \addplot[very thick, titanred, smooth] coordinates {
      (0,0) (0.3,6) (0.5,8) (1,8) (1.5,6) (2,4) (2.5,2)
      (2.7,3) (3,8) (3.5,10) (4,12) (4.5,10) (5,6)
      (5.2,4) (5.5,8) (6,14) (6.5,16) (7,10)
      (7.2,5) (7.5,3) (7.7,2) (8,1) (8.5,1) (9,0)
    };
    \addlegendentry{Враги на арене}

    % Intensity overlay
    \addplot[very thick, titanblue, dashed, smooth] coordinates {
      (0,1) (0.5,5) (1,6) (1.5,5) (2,4) (2.5,2)
      (3,5) (3.5,7) (4,9) (4.5,8) (5,5)
      (5.5,6) (6,10) (6.5,14) (7,12)
      (7.5,16) (8,20) (8.5,15) (9,5)
    };
    \addlegendentry{Интенсивность}

    % Breath markers
    \node[font=\tiny, titangreen, fill=white, inner sep=1pt]
      at (axis cs:2.5,3) {Breath};
    \node[font=\tiny, titangreen, fill=white, inner sep=1pt]
      at (axis cs:5,5) {Breath};
    \node[font=\tiny, titangreen, fill=white, inner sep=1pt]
      at (axis cs:7,8) {Breath};
  \end{axis}
\end{tikzpicture}

\subsection{Детальный план волн}

\subsubsection{Волна 1: Введение (0--2.5 мин)}

Цель: дать игроку освоиться в пространстве арены.

\begin{table}[H]
\centering
\begin{tabularx}{\textwidth}{@{}l c c X@{}}
\toprule
\textbf{Подволна} & \textbf{Состав} & \textbf{Зона} & \textbf{Цель} \\
\midrule
1a (0:00) & 4 Rusher + 2 Anchor & Ground & Знакомство с ареной, заставить двигаться \\
1b (0:30) & 2 Anchor & Mid (платформы) & Показать вертикальность арены \\
1c (1:00) & 2 Rusher + 1 Disruptor & Ground & Ввести AoE-давление, заставить летать \\
\bottomrule
\end{tabularx}
\end{table}

\subsubsection{Волна 2: Эскалация (2.5--5 мин)}

Цель: увеличить сложность, ввести Tank и Specialist.

\begin{table}[H]
\centering
\begin{tabularx}{\textwidth}{@{}l c c X@{}}
\toprule
\textbf{Подволна} & \textbf{Состав} & \textbf{Зона} & \textbf{Цель} \\
\midrule
2a (2:30) & 4R + 2A + 1T & Ground+Mid & Tank давит снизу, anchor-ы с платформ \\
2b (3:30) & 2D + 1Sp & Mid+High & Specialist на высоте, disruptors давят AoE \\
2c (4:00) & 3R & Ground & Подкрепление rusher-ов не даёт расслабиться \\
\bottomrule
\end{tabularx}
\end{table}

\subsubsection{Волна 3: Кульминация (5--7 мин)}

Цель: максимальное давление, все роли одновременно.

\begin{table}[H]
\centering
\begin{tabularx}{\textwidth}{@{}l c c X@{}}
\toprule
\textbf{Подволна} & \textbf{Состав} & \textbf{Зона} & \textbf{Цель} \\
\midrule
3a (5:00) & 4R + 3A + 2D & Все зоны & Максимальное кол-во, все направления \\
3b (5:30) & 1T + 1Sp & Ground+High & Tank внизу + specialist наверху = вилка \\
3c (6:00) & 4R + 2A & Mid & Финальная подволна, сброс перед боссом \\
\bottomrule
\end{tabularx}
\end{table}

\subsubsection{Босс-фаза (7--9 мин)}

Цель: финальное испытание, максимальный skill expression.

\begin{table}[H]
\centering
\begin{tabularx}{\textwidth}{@{}l c c X@{}}
\toprule
\textbf{Подволна} & \textbf{Состав} & \textbf{Зона} & \textbf{Цель} \\
\midrule
Boss (7:00) & 1 Boss & Центр & Boss использует всю арену; adds появляются по таймеру \\
Adds (7:30+) & 2R каждые 30с & Ground & Поддержание давления, ресурсная подпитка \\
\bottomrule
\end{tabularx}
\end{table}

\begin{polarity}{Интеграция полярности}
Каждая волна может менять \textbf{полярный заряд} арены:
\begin{itemize}
  \item Волна 1: нейтральный заряд~--- базовое состояние
  \item Волна 2: положительный заряд~--- одни платформы активны
  \item Волна 3: отрицательный заряд~--- другие платформы активны
  \item Босс: полярность меняется динамически, привязана к фазам босса
\end{itemize}
\end{polarity}


\section{Правила spawn-а на большой арене}

\begin{rulebox}{8 правил спавна в центральной арене}
\begin{enumerate}[leftmargin=*]
  \item \textbf{Никогда в прямом FoV игрока.} Враги не должны появляться
    из воздуха в центре экрана. Спавн за платформами, из порталов,
    из дверей, из-за стен
  \item \textbf{Минимум 20м от игрока.} При спавне враг должен быть
    достаточно далеко, чтобы игрок мог среагировать, но достаточно
    близко, чтобы быть заметным
  \item \textbf{Распределять по зонам.} Не кластерить всех врагов в
    одном углу. Минимум 2 зоны заняты врагами одновременно
  \item \textbf{Аудио-прелюдия.} Каждый спавн предваряется звуком:
    сирена, рёв, звук портала. Это даёт 0.5--1 секунду подготовки
  \item \textbf{Spawn on kill trigger.} Подволна 2b спавнится, когда
    убиты 50--70\% подволны 2a. Это создаёт непрерывность без перегрузки
  \item \textbf{Max concurrent enemies.} Одновременно на арене не
    больше 12--16 врагов. Больше = visual noise, меньше = скучно
  \item \textbf{Tank-и и Specialist-ы~--- по одному.} Максимум 1 Tank
    и 1 Specialist одновременно (до босс-фазы)
  \item \textbf{Не спавнить в Edge Zone.} Ресурсная зона у стен
    должна оставаться <<безопасной>> (относительно). Враги спавнятся
    в Kill Zone или Flank Zone
\end{enumerate}
\end{rulebox}


% =============================================================================
\chapter{Пространственные паттерны размещения}
% =============================================================================

\section{Паттерн <<Триада>> (3 врага)}

Минимальная осмысленная группа. Враги размещаются треугольником.

\begin{tikzpicture}[scale=0.8]
  \node[font=\bfseries\color{titandark}] at (5,6) {Паттерн <<Триада>>: три варианта};

  % Variant A: Even triangle
  \draw[thick, titandark!20, fill=titanlight] (0,0) rectangle (4,4);
  \fill[titangreen] (2,0.5) circle (0.25);
  \node[font=\tiny, titangreen] at (2,0) {Player};
  \fill[titanred] (1,3) circle (0.2);
  \fill[titanblue] (3,3) circle (0.2);
  \fill[titanorange] (2,2) circle (0.2);
  \node[font=\tiny, titanred] at (1,3.5) {R};
  \node[font=\tiny, titanblue] at (3,3.5) {A};
  \node[font=\tiny, titanorange] at (2,1.5) {D};
  \draw[dashed, titandark!30] (1,3) -- (3,3) -- (2,2) -- cycle;
  \node[font=\scriptsize\bfseries] at (2,-0.7) {A: R+A+D};
  \node[font=\tiny] at (2,-1.2) {Классическая};

  % Variant B: Depth triangle
  \draw[thick, titandark!20, fill=titanlight] (5.5,0) rectangle (9.5,4);
  \fill[titangreen] (7.5,0.5) circle (0.25);
  \fill[titanred] (6.5,1.5) circle (0.2);
  \fill[titanred] (8.5,1.5) circle (0.2);
  \fill[titanblue] (7.5,3.5) circle (0.2);
  \node[font=\tiny, titanred] at (6.5,1) {R};
  \node[font=\tiny, titanred] at (8.5,1) {R};
  \node[font=\tiny, titanblue] at (7.5,3) {A};
  \draw[dashed, titandark!30] (6.5,1.5) -- (8.5,1.5) -- (7.5,3.5) -- cycle;
  \node[font=\scriptsize\bfseries] at (7.5,-0.7) {B: 2R+A};
  \node[font=\tiny] at (7.5,-1.2) {Агрессивная};

  % Variant C: Sniper + rushers
  \draw[thick, titandark!20, fill=titanlight] (11,0) rectangle (15,4);
  \fill[titangreen] (13,0.5) circle (0.25);
  \fill[titanred] (12,2) circle (0.2);
  \fill[titanred] (14,2) circle (0.2);
  \fill[titancyan] (13,3.8) circle (0.2);
  \node[font=\tiny, titanred] at (12,1.5) {R};
  \node[font=\tiny, titanred] at (14,1.5) {R};
  \node[font=\tiny, titancyan] at (13,3.3) {Sp};
  \draw[dashed, titandark!30] (12,2) -- (14,2) -- (13,3.8) -- cycle;
  \node[font=\scriptsize\bfseries] at (13,-0.7) {C: 2R+Sp};
  \node[font=\tiny] at (13,-1.2) {Приоритетная};
\end{tikzpicture}

\vspace{5pt}
\begin{itemize}[leftmargin=*]
  \item \textbf{A (R+A+D):} Все три давления одновременно.
    Игрок не может стоять (R), прятаться (D) или игнорировать (A).
    Идеально для арены
  \item \textbf{B (2R+A):} Два rusher-а бегут с флангов,
    anchor стреляет из глубины. Игрок решает: убить близких
    или прорваться к anchor-у. Подходит для коридорных стычек
  \item \textbf{C (2R+Sp):} Rusher-ы не дают спокойно целиться
    в specialist-а. Создаёт дилемму приоритетов.
    Подходит для обеих частей арены
\end{itemize}


\section{Паттерн <<Тиски>> (2+2)}

Две группы зажимают игрока с разных направлений.

\begin{tikzpicture}[scale=0.8]
  \node[font=\bfseries\color{titandark}] at (7,6.5) {Паттерн <<Тиски>>};

  % Arena space
  \draw[thick, titandark!20, fill=titanlight] (0,0) rectangle (14,5);
  \fill[titangreen] (7,2.5) circle (0.3);
  \node[font=\scriptsize\bfseries, titangreen] at (7,1.8) {PLAYER};

  % Group A (left)
  \fill[titanred] (2,2) circle (0.2);
  \fill[titanred] (2.5,3) circle (0.2);
  \fill[titanblue] (1.5,3.5) circle (0.2);
  \draw[thick, titanred, -{Stealth}] (2.8,2.5) -- (5.5,2.5);
  \node[font=\scriptsize\bfseries, titanred] at (2,4.3) {Группа A};
  \node[font=\tiny, titandark] at (2,1.3) {2R + 1A};

  % Group B (right)
  \fill[titanred] (12,2.5) circle (0.2);
  \fill[titanorange] (11.5,3.5) circle (0.2);
  \fill[titanblue] (12.5,3) circle (0.2);
  \draw[thick, titanred, -{Stealth}] (11.2,2.8) -- (8.5,2.5);
  \node[font=\scriptsize\bfseries, titanred] at (12,4.3) {Группа B};
  \node[font=\tiny, titandark] at (12,1.3) {1R + 1A + 1D};

  % Escape routes
  \draw[very thick, titangreen, dashed, -{Stealth}] (7,2.5) -- (7,5.5);
  \node[font=\tiny, titangreen] at (7,5.8) {Побег вверх (полёт)};
  \draw[very thick, titangreen, dashed, -{Stealth}] (7,2.5) -- (7,-0.5);
  \node[font=\tiny, titangreen] at (7,-0.8) {Побег вниз};
\end{tikzpicture}

\begin{warning}{Правило выхода из тисков}
В паттерне <<тиски>> \textbf{ВСЕГДА} должен быть третий путь
побега (вверх, вниз, через портал). Если игрок зажат без
выхода~--- это не сложно, это \textbf{нечестно}. В 3D-арене
побег вверх/вниз~--- естественный и читаемый выход.
\end{warning}


\section{Паттерн <<Слои>> (вертикальное распределение)}

Специфичен для арены 100$\times$100$\times$50м с активным
полётом.

\begin{tikzpicture}[scale=0.75]
  \node[font=\bfseries\color{titandark}] at (7,10.5) {Паттерн <<Слои>>: вертикальное давление};

  % Arena cross-section
  \draw[very thick, titandark] (1,0) rectangle (13,9);

  % Layer 1: Ground
  \fill[titanblue!10] (1,0) rectangle (13,3);
  \fill[titangold] (4,1) circle (0.4);
  \node[font=\tiny, white] at (4,1) {Tank};
  \fill[titanred] (8,0.8) circle (0.2);
  \fill[titanred] (9,1.5) circle (0.2);
  \fill[titanred] (10,0.6) circle (0.2);
  \node[font=\scriptsize\bfseries, titanblue] at (7,2.5) {Layer 1: Ground};
  \node[font=\tiny, titandark] at (7,2) {Tank + Rushers};

  % Layer 2: Mid
  \fill[titanorange!8] (1,3) rectangle (13,6);
  \fill[titanblue!70] (3,4.5) circle (0.2);
  \fill[titanblue!70] (6,4) circle (0.2);
  \fill[titanblue!70] (10,5) circle (0.2);
  \fill[titanorange] (8,4.5) circle (0.2);
  \node[font=\scriptsize\bfseries, titanorange] at (7,5.5) {Layer 2: Mid};
  \node[font=\tiny, titandark] at (7,3.5) {Anchors + Disruptor};

  % Layer 3: High
  \fill[titanred!8] (1,6) rectangle (13,9);
  \fill[titancyan] (5,7.5) circle (0.2);
  \fill[titancyan] (9,8) circle (0.2);
  \node[font=\scriptsize\bfseries, titanred] at (7,8.5) {Layer 3: High};
  \node[font=\tiny, titandark] at (7,6.5) {Specialists};

  % Player in mid
  \fill[titangreen] (7,4.5) circle (0.3);
  \node[font=\tiny, white] at (7,4.5) {P};

  % Pressure arrows
  \draw[ultra thick, titangold!60, -{Stealth}] (4,1.5) -- (6.5,4);
  \draw[ultra thick, titancyan!60, -{Stealth}] (5,7.2) -- (6.5,5);
  \draw[ultra thick, titanblue!60, -{Stealth}] (3,4.5) -- (6.5,4.5);
  \draw[ultra thick, titanblue!60, -{Stealth}] (10,5) -- (7.5,4.7);
\end{tikzpicture}

\begin{principle}{Принцип слоёв}
Каждый вертикальный слой давит своим типом угрозы:
\begin{itemize}
  \item \textbf{Ground:} <<опуститься = умереть>> (Tank + Rushers)
  \item \textbf{Mid:} <<остановиться = умереть>> (Anchors + Disruptors)
  \item \textbf{High:} <<игнорировать = умереть>> (Specialists)
\end{itemize}
Игрок \textbf{вынужден} постоянно менять высоту, чтобы
избегать давления одного слоя, при этом попадая под
давление другого. Это создаёт непрерывный <<танец>>.
\end{principle}


\section{Паттерн <<Якорь и Рой>>}

Один сильный враг (Tank или мини-босс) привлекает внимание
на себя, пока рой мелких врагов (Rusher-ы) атакует с периферии.

\begin{tikzpicture}[scale=0.8]
  \draw[thick, titandark!20, fill=titanlight] (0,0) rectangle (10,6);
  \fill[titangreen] (5,1) circle (0.3);
  \node[font=\scriptsize, titangreen] at (5,0.4) {Player};

  % Tank in center
  \fill[titangold] (5,4) circle (0.5);
  \node[font=\small\bfseries, white] at (5,4) {TANK};

  % Rusher swarm
  \foreach \angle in {0,45,...,315} {
    \fill[titanred] ($(5,4)+(\angle:2)$) circle (0.15);
  }

  % Attention split
  \draw[ultra thick, titangold, -{Stealth}] (5,3.3) -- (5,1.8);
  \node[font=\tiny, titangold, right] at (5.3,2.5) {Главная угроза};
  \draw[thick, titanred, -{Stealth}] ($(5,4)+(200:2)$) -- ($(5,1)+(160:1.5)$);
  \draw[thick, titanred, -{Stealth}] ($(5,4)+(340:2)$) -- ($(5,1)+(20:1.5)$);
  \node[font=\tiny, titanred] at (8,2) {Периферийные угрозы};

  \node[font=\small, text=titandark, text width=9cm, align=center] at (5,-1) {
    Tank требует фокуса. Rusher-ы наказывают за тоннельное зрение.
    Игрок должен \textbf{переключать внимание} между главной и
    второстепенными угрозами.
  };
\end{tikzpicture}


% =============================================================================
\chapter{Правила честности и читаемости}
% =============================================================================

\section{Честность (Fairness)}

\begin{principle}{Золотое правило}
Каждая смерть игрока должна ощущаться как \textbf{его ошибка},
а не как подлость дизайнера. Если игрок не понимает, почему
он умер~--- энкаунтер спроектирован плохо.
\end{principle}

\begin{table}[H]
\centering
\caption{Честно vs Нечестно}
\begin{tabularx}{\textwidth}{@{}XX@{}}
\toprule
\textbf{Честно} & \textbf{Нечестно} \\
\midrule
Враг виден до атаки & Враг стреляет из невидимой позиции \\
Засада с аудио-подсказкой & Внезапный спавн без предупреждения \\
Сложная, но проходимая волна & 20 врагов спавнятся одновременно \\
Тиски с путём побега & Тиски без выхода \\
Sniper с лазерным прицелом & Instant-hit снайпер без индикации \\
Враг за укрытием~--- можно зафланкировать & Враг за неразрушимой стеной \\
\bottomrule
\end{tabularx}
\end{table}

\section{Читаемость (Readability)}

\subsection{Визуальная читаемость врагов}

\begin{rulebox}{Правила читаемости}
\begin{enumerate}[leftmargin=*]
  \item \textbf{Силуэт.} Каждая роль должна быть различима
    на расстоянии 10+ метров по силуэту
  \item \textbf{Цвет.} Роли должны иметь разные цветовые акценты
    (Rusher~--- красный, Anchor~--- синий, Disruptor~--- оранжевый,
    Tank~--- жёлтый, Specialist~--- бирюзовый)
  \item \textbf{Размер.} Tank > Anchor > Rusher по физическим размерам.
    Большой = опасный и медленный
  \item \textbf{Звук.} У каждой роли свои звуковые подсказки:
    топот (Rusher), жужжание (Disruptor), рёв (Tank),
    щелчок (Specialist/Sniper)
  \item \textbf{Телеграфирование атак.} Каждая атака имеет
    визуальный windup длительностью 0.3--0.8 секунды.
    Чем сильнее атака~--- тем длиннее windup
\end{enumerate}
\end{rulebox}

\subsection{Правило <<одного взгляда>>}

\begin{principle}{One Glance Rule}
Игрок должен понять состав энкаунтера \textbf{за один взгляд}
(0.5--1 секунда). Это значит:
\begin{itemize}
  \item Не более 3 типов врагов в одной группе (до Волны 3)
  \item Враги одного типа стоят \textbf{рядом друг с другом}
    (подсознательное группирование)
  \item Самый опасный враг визуально \textbf{выделяется}
    (размером, свечением, позицией)
\end{itemize}
\end{principle}


% =============================================================================
\chapter{Ресурсы и экономика энкаунтера}
% =============================================================================

\section{Ресурсная петля DOOM Eternal}

\begin{tikzpicture}[
  node distance=1.8cm,
  res/.style={
    rectangle, rounded corners=4pt,
    minimum width=2.8cm, minimum height=1.3cm,
    draw=#1, fill=#1!10,
    font=\small\bfseries, text=titandark, align=center
  },
  arr/.style={-{Stealth[length=5pt]}, very thick, #1}
]
  \node[res=titangreen] (hp) at (0,0) {HP\\{\scriptsize Glory Kill}};
  \node[res=titanblue] (ammo) at (5,0) {Ammo\\{\scriptsize Chainsaw}};
  \node[res=titanorange] (armor) at (2.5,3) {Armor\\{\scriptsize Flame Belch}};

  \draw[arr=titangreen] (ammo) -- (hp) node[midway, below, font=\tiny] {Стреляешь $\to$ ослабляешь $\to$ добиваешь};
  \draw[arr=titanblue] (hp) -- (armor) node[midway, left, font=\tiny, text width=2cm, align=center] {Живой $\to$ можешь атаковать};
  \draw[arr=titanorange] (armor) -- (ammo) node[midway, right, font=\tiny, text width=2cm, align=center] {Защищён $\to$ можешь рисковать};

  \node[font=\small, titandark, below=1cm of ammo, text width=10cm, align=center] {
    \textbf{Принцип:} враги одновременно \textbf{угроза} и \textbf{ресурс}.
    Убийство врага восполняет то, что он отнял. Это создаёт
    push-forward combat~--- агрессия вознаграждается.
  };
\end{tikzpicture}

\section{Размещение ресурсов на арене}

\begin{rulebox}{Правила размещения ресурсов}
\begin{enumerate}[leftmargin=*]
  \item \textbf{Ресурсы на периферии.} HP и аммо~--- в Edge Zone.
    Чтобы подобрать, нужно выбежать из Kill Zone (risk/reward)
  \item \textbf{Мелкие ресурсы~--- в центре.} Маленькие дропы~---
    из убитых врагов. Большие пикапы~--- у стен
  \item \textbf{Ресурсы указывают путь.} Цепочка маленьких ресурсов
    ведёт к большому~--- это и навигация, и награда за движение
  \item \textbf{Не ставить ресурсы рядом со спавном врагов.}
    Ресурс~--- это момент передышки, спавн~--- момент напряжения.
    Они не должны быть в одном месте
  \item \textbf{Критический ресурс~--- у Boss-а.} Если HP дропается
    при ударе по боссу~--- это мотивирует агрессию, а не кайтинг
\end{enumerate}
\end{rulebox}


% =============================================================================
\chapter{Чек-листы}
% =============================================================================

\section{Чек-лист: Коридорная стычка (Часть I)}

\begin{table}[H]
\centering
\begin{tabularx}{\textwidth}{@{}cX@{}}
\toprule
$\square$ & Группа = 2--4 врага, не более \\
$\square$ & Минимум 2 разные роли в группе (кроме первой стычки) \\
$\square$ & Враги привязаны к конкретному геометрическому элементу \\
$\square$ & Дистанция обнаружения: 15--30м (не ближе, не дальше) \\
$\square$ & Враги не спавнятся за спиной (кроме scripted засады) \\
$\square$ & Есть аудио-подсказка до визуального контакта \\
$\square$ & Между стычками есть 5--15 секунд передышки \\
$\square$ & Каждая следующая стычка сложнее предыдущей (эскалация) \\
$\square$ & Последняя стычка перед ареной~--- самая интенсивная \\
$\square$ & Есть vista (вид на арену) хотя бы через 1 окно \\
\bottomrule
\end{tabularx}
\end{table}

\section{Чек-лист: Волна на арене (Часть II)}

\begin{table}[H]
\centering
\begin{tabularx}{\textwidth}{@{}cX@{}}
\toprule
$\square$ & Макс. 12--16 врагов одновременно \\
$\square$ & Макс. 3 типа ролей одновременно (до Волны 3) \\
$\square$ & Враги распределены по минимум 2 вертикальным зонам \\
$\square$ & Есть Rusher-ы для давления на позицию \\
$\square$ & Есть Anchor-ы для давления на укрытие \\
$\square$ & Есть хотя бы один Disruptor ИЛИ Specialist \\
$\square$ & Нет <<безопасной высоты>> (каждый слой под давлением) \\
$\square$ & Ресурсы доступны на периферии арены \\
$\square$ & Спавн не в FoV игрока \\
$\square$ & Аудио-прелюдия перед каждым спавном \\
$\square$ & Есть breath (1--3 сек тишины) между волнами \\
$\square$ & Bosс-фаза начинается после финальной <<тишины>> \\
\bottomrule
\end{tabularx}
\end{table}

\section{Чек-лист: Комбинация врагов}

\begin{table}[H]
\centering
\begin{tabularx}{\textwidth}{@{}cX@{}}
\toprule
$\square$ & Нет однородных групп (минимум 2 роли, кроме обучения) \\
$\square$ & Rusher + Anchor = стандартная база каждой волны \\
$\square$ & Каждая комбинация создаёт <<конфликт тактик>> \\
$\square$ & Нет единственной безопасной тактики \\
$\square$ & Tank и Specialist~--- максимум по 1 одновременно \\
$\square$ & Каждый враг читаем за 1 секунду (силуэт + цвет + звук) \\
$\square$ & Каждая атака имеет windup 0.3--0.8 сек \\
$\square$ & Каждая смерть игрока~--- его ошибка, не подлость дизайнера \\
\bottomrule
\end{tabularx}
\end{table}


% =============================================================================
\chapter{Итерация и плейтестинг}
% =============================================================================

\section{Что тестировать}

\begin{table}[H]
\centering
\caption{Метрики плейтеста для энкаунтеров}
\begin{tabularx}{\textwidth}{@{}lXX@{}}
\toprule
\textbf{Метрика} & \textbf{Хорошо} & \textbf{Плохо} \\
\midrule
Смертей на стычку (Ч.I) & 0--1 (Normal) & 3+ = слишком сложно \\
Смертей на волну (Ч.II) & 0--2 (Normal) & 4+ = проблема с fairness \\
Время на стычку (Ч.I) & 15--45 сек & 60+ сек = скучно \\
Время на волну (Ч.II) & 1.5--3 мин & 4+ мин = затянуто \\
\% использования оружий & Равномерно & Одно оружие = слабая палитра врагов \\
Тепловая карта смертей & Распределена & Кластер = проблемная точка \\
Движение игрока & По всей арене & Сидит в углу = плохие стимулы \\
Вертикальность & Использует все слои & Только ground = плохая вертикаль \\
\bottomrule
\end{tabularx}
\end{table}

\section{Процесс итерации}

\begin{enumerate}[leftmargin=*]
  \item \textbf{Грейбокс + placeholder враги.} Расставьте врагов
    базовых типов, запустите, пройдите сами 5 раз
  \item \textbf{Запишите видео.} Каждый проход~--- запись.
    Отмечайте, где умерли, где было скучно, где было <<вау>>
  \item \textbf{Тепловая карта.} Где игрок проводит больше всего
    времени? Если он стоит в одном месте~--- нужен Disruptor
    в этой зоне
  \item \textbf{Внешний плейтест.} Дайте пройти 3--5 людям без
    подсказок. Спросите: <<где было непонятно?>> и
    <<где было нечестно?>>
  \item \textbf{Итерация.} По результатам: сдвиньте спавн на 5м,
    уберите одного врага, добавьте укрытие. Маленькие изменения~---
    большой эффект
\end{enumerate}

\begin{tip}{Правило <<трёх проходов>>}
Каждый энкаунтер нужно пройти минимум:
\begin{itemize}
  \item \textbf{Агрессивно}~--- бежать вперёд, стрелять на бегу
  \item \textbf{Осторожно}~--- прятаться, стрелять из укрытия
  \item \textbf{Стильно}~--- использовать все механики, combo, полёт
\end{itemize}
Если один из проходов \textbf{невозможен}~--- геометрия или палитра
врагов ограничивает игрока слишком сильно.
Если один из проходов \textbf{явно лучший}~--- нужно добавить врагов,
которые наказывают эту тактику.
\end{tip}

\vspace{1cm}

\begin{center}
\begin{tikzpicture}
  \fill[titandark, rounded corners=4pt] (0,0) rectangle (12,2);
  \node[text=white, font=\Large\bfseries] at (6,1.3) {POLARITY};
  \node[text=titanlight, font=\small] at (6,0.6) {Encounter Design Guide --- Internal Document --- 2026};
\end{tikzpicture}
\end{center}

\end{document}
